% TOP_PROBLEM: {"id": 14, "title": "Hadwiger Conjecture", "category_id": 3, "category": {"id": 3, "name": "graph_theory", "display_name": "Graph Theory", "description": "Problems involving graphs, networks, and their properties.", "slug": "graph-theory", "order_index": 3, "created_at": "2026-07-31T15:26:25.671Z"}, "status": "open"}
% FIRST_POSTED: 2026-10-06
% ATTEMPT_STATUS: solved
% OPENAI_CLAIM: {"schema_version":1,"source_repo":"https://github.com/openai/math","source_commit":"adc7f1241b42e322a6451854ab7e4b4c146bf78a","release_date":"2026-10-06","match":"full","resolution":"disproved","independently_reviewed":false,"second_pass":"agreed","families":[{"family":"157","title":"Graph coloring, clique minors, and Colin de Verdière invariants","subject":"Combinatorics","manuscripts":[{"title":"A counterexample to Hadwiger's conjecture","pdf_path":"preprints/A-counterexample-to-Hadwigers-conjecture-September-23-2026/paper.pdf","readme_path":"preprints/A-counterexample-to-Hadwigers-conjecture-September-23-2026/README.md","date":"2026-09-23","theorem_ref":"thm:main (sections/01-introduction.tex:20); cor:hadwiger (sections/01-introduction.tex:48)","statement_sources":[{"path":"preprints/A-counterexample-to-Hadwigers-conjecture-September-23-2026/build/sections/01-introduction.tex","line":20,"label":"thm:main"},{"path":"preprints/A-counterexample-to-Hadwigers-conjecture-September-23-2026/build/sections/01-introduction.tex","line":48,"label":"cor:hadwiger"}]},{"title":"A counterexample to the Colin de Verdière chromatic conjecture","pdf_path":"preprints/A-counterexample-to-the-Colin-de-Verdiere-chromatic-conjecture-September-23-2026/paper.pdf","readme_path":"preprints/A-counterexample-to-the-Colin-de-Verdiere-chromatic-conjecture-September-23-2026/README.md","date":"2026-09-23","theorem_ref":"thm:main (sections/01-introduction.tex:42); thm:base-graphs (sections/01-introduction.tex:147); cor:fractional-spectral (sections/01-introduction.tex:221); cor:list-spectral (sections/01-introduction.tex:262)","statement_sources":[{"path":"preprints/A-counterexample-to-the-Colin-de-Verdiere-chromatic-conjecture-September-23-2026/build/sections/01-introduction.tex","line":42,"label":"thm:main"},{"path":"preprints/A-counterexample-to-the-Colin-de-Verdiere-chromatic-conjecture-September-23-2026/build/sections/01-introduction.tex","line":147,"label":"thm:base-graphs"},{"path":"preprints/A-counterexample-to-the-Colin-de-Verdiere-chromatic-conjecture-September-23-2026/build/sections/01-introduction.tex","line":221,"label":"cor:fractional-spectral"},{"path":"preprints/A-counterexample-to-the-Colin-de-Verdiere-chromatic-conjecture-September-23-2026/build/sections/01-introduction.tex","line":262,"label":"cor:list-spectral"}]},{"title":"A linear list-coloring bound in terms of the Hadwiger number","pdf_path":"preprints/A-linear-list-coloring-bound-in-terms-of-the-Hadwiger-number-September-23-2026/paper.pdf","readme_path":"preprints/A-linear-list-coloring-bound-in-terms-of-the-Hadwiger-number-September-23-2026/README.md","date":"2026-09-23","theorem_ref":"thm:main (sections/01-introduction.tex:25); cor:ordinary (sections/01-introduction.tex:38)","statement_sources":[{"path":"preprints/A-linear-list-coloring-bound-in-terms-of-the-Hadwiger-number-September-23-2026/build/sections/01-introduction.tex","line":25,"label":"thm:main"},{"path":"preprints/A-linear-list-coloring-bound-in-terms-of-the-Hadwiger-number-September-23-2026/build/sections/01-introduction.tex","line":38,"label":"cor:ordinary"}]}],"lean":{"doc_path":"lean/docs/157.md","comparators":["lean/ComparatorChallenges/ListHadwiger.json","lean/ComparatorChallenges/ListHadwiger.lean"],"declarations":[],"covers_main_theorem":false},"reasoning_trace":null}]}
\documentclass[11pt]{article}
\usepackage[margin=25mm]{geometry}
\usepackage{fontspec}
\setmainfont{Latin Modern Roman}
\usepackage{amsmath,amssymb,mathtools}
\usepackage{unicode-math}
\setmathfont{Latin Modern Math}
\usepackage{enumitem}
\usepackage{xurl,hyperref,bookmark}
\hypersetup{colorlinks=true,urlcolor=blue}
\setcounter{secnumdepth}{0}
\setlength{\parindent}{0pt}
\setlength{\parskip}{5pt}
\setlength{\emergencystretch}{3em}
\begin{document}
\section*{Hadwiger Conjecture}

\subsection{OpenAI's claim}
\textit{Chromatic number beyond clique-minor size}: The counterexample paper claims arbitrarily large finite graphs with $\chi_f(G)>h(G)$, hence also $\chi(G)>h(G)$; no independent vertex set in an example contains three vertices.

\textit{Colin de Verdière chromatic obstruction}: The manuscript claims finite graphs whose fractional chromatic number is greater than $\mu(G)+1$, and therefore whose ordinary chromatic number violates the same bound; such examples exist at arbitrarily large orders.

\textit{Linear list-coloring in clique-minor size}: The paper claims a universal constant $C$ with $\chi_{\mathrm{list}}(G)\le C h(G)$ for every finite nonempty graph; the factor is independent of the graph and its order.

\subsection{Mathematical statement of the claimed result}
\paragraph{Mathematical claim 1.} There are finite simple graphs $G$ of arbitrarily large order $m$ with independence number $\alpha(G)\le2$ and maximum size of a pairwise-touching matching less than $m/100$. The resulting corollary states \[h(G)<26m/75+2/3<m/2\le\chi_f(G)\le\chi(G),\] where $h(G)=\max\{t:K_t\text{ is a minor of }G\}$ and $\chi_f,\chi$ are fractional and ordinary chromatic numbers. Thus the exact inequalities $\chi\le h$ and $\chi_f\le h$ fail; no conclusion about replacing $h$ by a fractional minor parameter is supplied.

\paragraph{Mathematical claim 2.} The claimed counterexamples are finite simple graphs of arbitrarily large order $m$ with \[\alpha(G)\le2,\qquad\mu(G)+1<m/2\le\chi(G),\] where $\mu(G)$ is the Colin de Verdi\`ere spectral graph invariant. The fractional-coloring corollary strengthens the obstruction to \[\chi_f(G)\ge m/\alpha(G)\ge m/2>\kappa(G)+1\ge\mu(G)+1\ge h(G).\] Here $\kappa$ is its strong-Arnold-property-free spectral relaxation, and $h$ is the largest complete-minor order. In particular both $\chi\le\mu+1$ and $\chi_f\le h$ are refuted. A further corollary gives $\chi_{\mathrm{list}}(G)\le C(\mu(G)+1)$ with an absolute $C$.

\paragraph{Mathematical claim 3.} For every finite nonempty simple graph $G$ there is an absolute integer $C\ge1$, independent of $G$, such that \[\chi_{\mathrm{list}}(G)\le C h(G),\qquad h(G)=\max\{t:K_t\text{ is a minor of }G\}.\] The list chromatic number is the least $k$ for which every assignment of $k$ available colors per vertex permits a proper coloring. Consequently $\chi(G)\le Ch(G)$. The constant-factor bound does not imply the conjectured coefficient $1$.

\subsection{Outline of the argument}
\textit{Chromatic number beyond clique-minor size}: The graph construction uses Boolean moment representations and paired frames to control compatible vertex distributions. Entropy projection, fingerprints for independent sets, and a capacitated-cut argument transfer the distributional constraints to finite graphs. The resulting obstruction limits clique-minor size, while the small independence number gives a larger fractional chromatic number.

\textit{Colin de Verdière chromatic obstruction}: Allowed channel images are thinned while preserving the primal distribution, supplying enough compatible pairs for the conflict estimates. A Boolean moment argument restricts low-rank matrix representations. Entropy and finite sampling then produce graphs for which every well-signed matrix with a single negative eigenvalue has too much rank to support the proposed chromatic inequality.

\textit{Linear list-coloring in clique-minor size}: For small minor-free graphs, a random global partition of colors gives each vertex many eligible bins. Stable-set sampling repeatedly reduces fractional chromatic complexity among eligible vertices. If the reduction failed, separator and connectivity arguments would assemble disjoint clique models into the excluded minor; the resulting small-graph bound feeds the global list-coloring argument.

\subsection{Scope relative to this problem}
\textit{Chromatic number beyond clique-minor size}: The claimed graphs have h(G)<chi\_f(G)<=chi(G), directly violating Hadwiger’s ordinary minor-coloring implication.

\textit{Colin de Verdière chromatic obstruction}: The claimed graphs have h(G)<chi\_f(G)<=chi(G), directly violating Hadwiger’s ordinary minor-coloring implication.

\textit{Linear list-coloring in clique-minor size}: The claimed list-coloring estimate implies chi(G)<=C h(G), a weaker constant-factor bound than Hadwiger’s required factor one.

The claimed counterexample lies within the question’s hypotheses and contradicts its universal conclusion.

\subsection{Formal verification}
Family 157 describes Lean coverage for the list-coloring upper bound with a universal multiple of clique-minor size; neither chromatic counterexample is part of this selection.

\textit{Chromatic number beyond clique-minor size}: No manuscript-specific Comparator is listed in the pinned inventory.

The documented selection does not cover the entire manuscript assertion relevant here.

\textit{Colin de Verdière chromatic obstruction}: No manuscript-specific Comparator is listed in the pinned inventory.

The documented selection does not cover the entire manuscript assertion relevant here.

\textit{Linear list-coloring in clique-minor size}: The manuscript links \texttt{ListHadwiger.lean}.

The documented selection includes this manuscript’s main assertion; that assertion retains the restrictions stated above.

This site has not run the Lean development. These links describe OpenAI’s selected formal statements and should be read with their stated scope.

\subsection{Status caveat}
OpenAI reports that an internal model produced the manuscripts. Their arguments have not been independently verified as proofs by this site. Any portion outside the documented formal scope may contain mistakes, and the catalogue label reports the announced mathematical claim.

\subsection{Sources}
\begin{itemize}
\item OpenAI, \emph{A counterexample to Hadwiger's conjecture}, 2026-09-23; citation key \texttt{OAI:A-counterexample-to-Hadwigers-conjecture-September-23-2026}. \href{https://github.com/openai/math/blob/adc7f1241b42e322a6451854ab7e4b4c146bf78a/preprints/A-counterexample-to-Hadwigers-conjecture-September-23-2026/paper.pdf}{PDF (GitHub)}; \href{https://github.com/openai/math/raw/adc7f1241b42e322a6451854ab7e4b4c146bf78a/preprints/A-counterexample-to-Hadwigers-conjecture-September-23-2026/paper.pdf}{PDF (direct)}; \href{https://github.com/openai/math/blob/adc7f1241b42e322a6451854ab7e4b4c146bf78a/preprints/A-counterexample-to-Hadwigers-conjecture-September-23-2026/README.md}{Manuscript page and citation}; \href{https://github.com/openai/math/blob/main/preprints/A-counterexample-to-Hadwigers-conjecture-September-23-2026/paper.pdf}{Latest version}; Statement source: \href{https://github.com/openai/math/blob/adc7f1241b42e322a6451854ab7e4b4c146bf78a/preprints/A-counterexample-to-Hadwigers-conjecture-September-23-2026/build/sections/01-introduction.tex\#L20}{thm:main}; Statement source: \href{https://github.com/openai/math/blob/adc7f1241b42e322a6451854ab7e4b4c146bf78a/preprints/A-counterexample-to-Hadwigers-conjecture-September-23-2026/build/sections/01-introduction.tex\#L48}{cor:hadwiger}.
\item OpenAI, \emph{A counterexample to the Colin de Verdière chromatic conjecture}, 2026-09-23; citation key \texttt{OAI:A-counterexample-to-the-Colin-de-Verdiere-chromatic-conjecture-September-23-2026}. \href{https://github.com/openai/math/blob/adc7f1241b42e322a6451854ab7e4b4c146bf78a/preprints/A-counterexample-to-the-Colin-de-Verdiere-chromatic-conjecture-September-23-2026/paper.pdf}{PDF (GitHub)}; \href{https://github.com/openai/math/raw/adc7f1241b42e322a6451854ab7e4b4c146bf78a/preprints/A-counterexample-to-the-Colin-de-Verdiere-chromatic-conjecture-September-23-2026/paper.pdf}{PDF (direct)}; \href{https://github.com/openai/math/blob/adc7f1241b42e322a6451854ab7e4b4c146bf78a/preprints/A-counterexample-to-the-Colin-de-Verdiere-chromatic-conjecture-September-23-2026/README.md}{Manuscript page and citation}; \href{https://github.com/openai/math/blob/main/preprints/A-counterexample-to-the-Colin-de-Verdiere-chromatic-conjecture-September-23-2026/paper.pdf}{Latest version}; Statement source: \href{https://github.com/openai/math/blob/adc7f1241b42e322a6451854ab7e4b4c146bf78a/preprints/A-counterexample-to-the-Colin-de-Verdiere-chromatic-conjecture-September-23-2026/build/sections/01-introduction.tex\#L42}{thm:main}; Statement source: \href{https://github.com/openai/math/blob/adc7f1241b42e322a6451854ab7e4b4c146bf78a/preprints/A-counterexample-to-the-Colin-de-Verdiere-chromatic-conjecture-September-23-2026/build/sections/01-introduction.tex\#L147}{thm:base-graphs}; Statement source: \href{https://github.com/openai/math/blob/adc7f1241b42e322a6451854ab7e4b4c146bf78a/preprints/A-counterexample-to-the-Colin-de-Verdiere-chromatic-conjecture-September-23-2026/build/sections/01-introduction.tex\#L221}{cor:fractional-spectral}; Statement source: \href{https://github.com/openai/math/blob/adc7f1241b42e322a6451854ab7e4b4c146bf78a/preprints/A-counterexample-to-the-Colin-de-Verdiere-chromatic-conjecture-September-23-2026/build/sections/01-introduction.tex\#L262}{cor:list-spectral}.
\item OpenAI, \emph{A linear list-coloring bound in terms of the Hadwiger number}, 2026-09-23; citation key \texttt{OAI:A-linear-list-coloring-bound-in-terms-of-the-Hadwiger-number-September-23-2026}. \href{https://github.com/openai/math/blob/adc7f1241b42e322a6451854ab7e4b4c146bf78a/preprints/A-linear-list-coloring-bound-in-terms-of-the-Hadwiger-number-September-23-2026/paper.pdf}{PDF (GitHub)}; \href{https://github.com/openai/math/raw/adc7f1241b42e322a6451854ab7e4b4c146bf78a/preprints/A-linear-list-coloring-bound-in-terms-of-the-Hadwiger-number-September-23-2026/paper.pdf}{PDF (direct)}; \href{https://github.com/openai/math/blob/adc7f1241b42e322a6451854ab7e4b4c146bf78a/preprints/A-linear-list-coloring-bound-in-terms-of-the-Hadwiger-number-September-23-2026/README.md}{Manuscript page and citation}; \href{https://github.com/openai/math/blob/main/preprints/A-linear-list-coloring-bound-in-terms-of-the-Hadwiger-number-September-23-2026/paper.pdf}{Latest version}; Statement source: \href{https://github.com/openai/math/blob/adc7f1241b42e322a6451854ab7e4b4c146bf78a/preprints/A-linear-list-coloring-bound-in-terms-of-the-Hadwiger-number-September-23-2026/build/sections/01-introduction.tex\#L25}{thm:main}; Statement source: \href{https://github.com/openai/math/blob/adc7f1241b42e322a6451854ab7e4b4c146bf78a/preprints/A-linear-list-coloring-bound-in-terms-of-the-Hadwiger-number-September-23-2026/build/sections/01-introduction.tex\#L38}{cor:ordinary}.
\item Family 157: \href{https://github.com/openai/math/blob/adc7f1241b42e322a6451854ab7e4b4c146bf78a/lean/docs/157.md}{Lean scope}.
\item \href{https://github.com/openai/math/blob/adc7f1241b42e322a6451854ab7e4b4c146bf78a/lean/ComparatorChallenges/ListHadwiger.json}{Comparator ListHadwiger.json}.
\item \href{https://github.com/openai/math/blob/adc7f1241b42e322a6451854ab7e4b4c146bf78a/lean/ComparatorChallenges/ListHadwiger.lean}{Comparator ListHadwiger.lean}.
\item \href{https://github.com/openai/math/blob/adc7f1241b42e322a6451854ab7e4b4c146bf78a/CONTENTS.md}{OpenAI catalogue entry}.
\end{itemize}
\end{document}
